\section{Kill Chain 与 Value Chain：共享软件骨架，不共享伦理结论}

Sankar 把政府 kill chain 与商业 value chain 视为结构相似的 decision chain：组织设定目标，建立 operational picture，分配资源，执行 workflow，再观察结果。技术平台可以复用 entity model、constraint engine、workflow、audit 和 feedback；政府使用武力、企业配置库存或资本的权限基础并不相同，不能因软件结构相似就抹平伦理差异。

\begin{figure}[H]
\centering
\includegraphics[width=0.94\textwidth]{images/09-common-decision-chain.png}
\caption{政府和商业系统共享 sensing、modeling、allocation 与 feedback 骨架，但保留不同 authority。}
\figsource{依据字幕 25:00--30:00 重绘，`images/09-common-decision-chain.png`。}
\end{figure}

\begin{knowledgebox}{读图：可复用的是 primitive，不是最终目标}
Shared software 可以提供 ontology、版本化规则、resource constraint、workflow 和 audit trail。政府目标受宪法、法律和指挥授权约束，商业目标受合同、监管、员工和客户权利约束。平台应让这些规则可配置且可审计，而不是把某个行业的默认目标偷偷带到另一个行业。
\end{knowledgebox}

\begin{table}[H]
\centering
\begin{tabular}{L{0.18\textwidth}L{0.34\textwidth}L{0.38\textwidth}}
\toprule
Primitive & 通用能力 & 必须外部定义的内容 \\
\midrule
Ontology & 统一 entity、relation、action & 哪些分类合法、谁能查看 \\
Constraint engine & 检查资源、政策和兼容性 & 规则来源与优先级 \\
Workflow & 协调人和系统步骤 & 哪些步骤必须人工批准 \\
Optimization & 在约束下选择方案 & Objective 与不可交易底线 \\
Audit trail & 记录输入、决策与执行 & 保存期限、监督者与纠错权 \\
Feedback & 用结果更新策略 & 哪些 outcome 算成功或伤害 \\
\bottomrule
\end{tabular}
\caption{技术平台提供 decision primitive，合法目标与权利边界必须由组织和制度定义。}
\end{table}

\subsection{本章小结}

政府和商业可以共享决策基础设施，但不能共享未经审查的 objective。Ontology、workflow、constraint 和 audit 是可复用 primitive；authority、rights 和 harm threshold 需要领域制度提供。这个区分也适用于 AI：模型供给可以通用，价值产生依赖具体需求系统。

